\documentclass[11pt,a4paper]{article}
\usepackage[margin=0.85in]{geometry}
\usepackage{graphicx}
\usepackage{booktabs}
\usepackage{hyperref}
\hypersetup{
  colorlinks=true,
  linkcolor=black,
  citecolor=black,
  urlcolor=blue,
  pdfborder={0 0 0}
}
\usepackage{amsmath,amssymb}
\usepackage{caption}
\usepackage{float}
\usepackage{placeins}
\usepackage{setspace}
\usepackage{microtype}
\usepackage{fancyhdr}
\usepackage{titlesec}
\usepackage{enumitem}
\setstretch{1.12}
\graphicspath{{outputs/}}
\setlist[itemize]{leftmargin=1.4em,itemsep=0.25em,topsep=0.35em}

\captionsetup{font=small,labelfont=bf,skip=4pt}
\setlength{\textfloatsep}{8pt plus 2pt minus 2pt}
\setlength{\floatsep}{6pt plus 2pt minus 2pt}
\setlength{\intextsep}{8pt plus 2pt minus 2pt}
\titleformat{\section}{\large\bfseries}{\thesection}{0.6em}{}
\titlespacing*{\section}{0pt}{1.1ex plus 0.3ex}{0.6ex}
\titleformat{\paragraph}[runin]{\bfseries}{}{0pt}{}[.~]
\titlespacing*{\paragraph}{0pt}{0.8ex}{0.5em}

\newcommand{\safeimg}[2][0.92\linewidth]{%
  \includegraphics[width=#1,height=0.38\textheight,keepaspectratio]{#2}%
}

\pagestyle{fancy}
\fancyhf{}
\fancyhead[L]{\footnotesize CISC3024 AI Assignment \#2}
\fancyhead[R]{\footnotesize CHE CHI HIN, Angus (UC325182)}
\fancyfoot[C]{\thepage}
\renewcommand{\headrulewidth}{0.3pt}

\title{\textbf{AI Assignment \#2}\\[0.25em]
\large Variational Bayesian Last Layers for\\
Clothing Image Classification}
\author{CHE CHI HIN, Angus\\Student ID: UC325182\\CISC3024 Pattern Recognition}
\date{October 2026}

\begin{document}
\maketitle
\thispagestyle{fancy}

\begin{abstract}
\textbf{Thesis:} a discriminative variational Bayesian last layer (D-VBLL; Harrison, Willes, and Snoek, ICLR 2024) turns the final linear classifier into a Gaussian posterior over weights.
On Fashion-MNIST, with a shared small CNN, D-VBLL matches a softmax MAP baseline on accuracy and negative log-likelihood and lowers expected calibration error.
The same posterior does not by itself detect MNIST digits as out-of-distribution, and neither model notices low-contrast shift, because the learned weight variance stays far below the prior and the predictive distribution stays close to a point-estimate softmax.
Search, implementation, experiments, and this report were done by a Cursor agent under a standing instruction that I would not hand-write code.
\end{abstract}

\section{How I Asked AI Tools to Find the Algorithm}
The Moodle brief asks for a \emph{recent} classification or recognition algorithm based on probability or Bayes theory, for computer vision, with the search, code, experiments, and report done by an AI tool.
I used Cursor only. I pasted the brief, pointed at my Assignment~1 repository, and gave the agent authority to choose the method.
Assignment~1 had already used FasterNet, a CVPR 2023 convolution designed for speed.
That model is not a probabilistic classifier, so it could not be reused as the answer to this brief.

The agent searched recent Bayesian deep-learning papers and kept three candidates that are actually about classification uncertainty:
\begin{itemize}
  \item \textbf{Variational Bayesian Last Layers} (Harrison, Willes, and Snoek, ICLR 2024), a sampling-free variational posterior on the last layer.
  \item \textbf{FSP-Laplace} (Cinquin et al., NeurIPS 2024), a Laplace approximation with a function-space Gaussian-process prior.
  \item \textbf{Last Layer Empirical Bayes} (Villecroze et al., 2025), a normalizing-flow prior on the last layer.
\end{itemize}
The search settled on D-VBLL, the discriminative classifier in that ICLR paper.
The training objective is a closed-form lower bound (their Theorem~2), so the agent can check the formula on a toy batch before any long run.
FSP-Laplace needs function-space machinery that is hard to verify on a laptop, and the empirical-Bayes flow adds another approximate model on top of the classifier.
The application is clothing recognition on Fashion-MNIST, which is a different task from the digit robustness study in Assignment~1.
The diary of prompts, rejected alternatives, and checks is \texttt{process\_log.md}.

\section{Algorithm Description}
\label{sec:algo}
\paragraph{Model.}
A network with weights $\theta$ maps an image $x$ to a feature vector $\phi(x,\theta)\in\mathbb{R}^{D}$.
Class logits are a noisy linear readout. For a last-layer matrix $W$ whose $k$-th row is $w_k$,
\begin{equation}
p(y\mid x,W,\theta)=\mathrm{softmax}(z),\qquad
z = W\phi(x,\theta)+\varepsilon,
\label{eq:disc}
\end{equation}
with $\varepsilon$ zero-mean Gaussian noise.
This is Bayesian logistic regression on learned features.
The prior on each row is isotropic Gaussian, $p(w_k)=\mathcal{N}(0, s^{2}I)$, with the authors' fan-in scaling $s^{2}=2/D$.
Exact softmax marginalization is intractable, so training uses a variational posterior that factorizes across classes,
$q(w_k)=\mathcal{N}(m_k,\mathrm{diag}(v_k))$.
Diagonal covariances drop cross-class correlation and make every loss term linear in $D$.

\paragraph{Training bound.}
Stochastic variational inference~\cite{hoffman2013svi} maximizes a lower bound on the log marginal likelihood.
For one example with label $y$, Theorem~2 replaces the expected log-softmax by
\begin{equation}
m_y^{\top}\phi
- \log\sum_{k}\exp\!\Big(m_k^{\top}\phi
+ \tfrac12\big(\phi^{\top}\mathrm{diag}(v_k)\phi+\sigma_k^{2}\big)\Big).
\label{eq:bound}
\end{equation}
The extra $\tfrac12\phi^{\top}\mathrm{diag}(v_k)\phi$ term is the Jensen gap: uncertain weights inflate every logit before the log-sum-exp, which stops the bound from behaving like an over-confident point estimate.
The mini-batch loss I minimize is the negative of that bound, plus $(1/T)\,\mathrm{KL}(q(W)\,\|\,p(W))$, minus a $(1/T)$ inverse-Wishart-style penalty on the noise variances $\sigma_k^{2}$.
$T$ is the number of training images.
Feature weights $\theta$ get an ordinary Gaussian prior, implemented as AdamW weight decay.
The last-layer parameters are not weight-decayed, because the KL term is already their regularizer.
This is the paper's ``full training'' objective, not a post-hoc Laplace approximation and not Bayes-by-Backprop on every layer~\cite{blundell2015weight}.

\paragraph{Prediction.}
The posterior predictive probability is the expectation of the softmax under $q(W)$.
That expectation has no closed form for the discriminative model, so I draw logit samples
\begin{equation}
z^{(s)}\sim\mathcal{N}\!\big(m^{\top}\phi,\;\mathrm{diag}(\phi^{\top}S\phi+\sigma^{2})\big)
\label{eq:pred}
\end{equation}
and average $\mathrm{softmax}(z^{(s)})$.
Test metrics use 32 samples.
Out-of-distribution scores use the maximum predictive probability~\cite{hendrycks2017baseline}.

\paragraph{What is Bayesian here.}
The quantity being inferred is the last-layer weight, not the image class alone.
The class probability reported at test time is a marginal probability: uncertainty in $W$ is integrated out.
A softmax MAP baseline uses the same features but a single point $W$.

\begin{figure}[H]
\centering
\safeimg[0.98\linewidth]{architecture.png}
\caption{Student-scale D-VBLL. The MAP baseline replaces the variational last layer with one linear layer and a softmax.}
\label{fig:arch}
\end{figure}

\section{How AI Implemented the Algorithm}
The Cursor agent wrote the repository.
I did not hand-write or hand-edit the Python.
The layout is:
\begin{itemize}
  \item \texttt{models/vbll.py} --- diagonal D-VBLL: Theorem~2 bound, closed-form KL, noise penalty, and Monte Carlo predictive probabilities.
  \item \texttt{models/backbone.py} --- shared CNN. Three conv stages, batch-norm, global average pool, 128-d features. About $1.58\times 10^{5}$ parameters.
  \item \texttt{models/classifier.py} --- \texttt{map} (cross-entropy) or \texttt{vbll} (ELBO). Weight decay applies to the backbone only when the head is Bayesian.
  \item \texttt{data\_utils.py} --- Fashion-MNIST / MNIST tensors, shared augmentation, pixel-space corruptions.
  \item \texttt{train.py}, \texttt{evaluate.py}, \texttt{metrics.py}, \texttt{figures.py}, \texttt{run\_all.py}, \texttt{run\_ablation.py}.
\end{itemize}
Before training, \texttt{self\_check()} builds an $8\times 4$ random batch and checks three facts: the loss is finite and back-propagates; the bound matches a hand-expanded log-sum-exp; multiplying the features by four increases logit variance.
That check is the implementation test.
It does not use Fashion-MNIST accuracy as a substitute for the formula.

Practical settings follow the paper and the authors' diagonal parameterization rather than a new variant.
The prior variance is $2/D$.
Weight standard deviations are initialized near $1/D$.
The KL is scaled by $1/T$ in the main runs.
I did not import the \texttt{vbll} library; the module is a direct implementation of the bound above so the loss can be read against Equation~\eqref{eq:bound}.

\section{Experiment Settings and Results}
\paragraph{Task.}
In-distribution data are Fashion-MNIST clothing images, $28\times 28$ greyscale, 10 classes~\cite{xiao2017fashion}.
The official training set is split once, with seed 0, into 54{,}000 training images and 6{,}000 validation images.
The 10{,}000 official test images are used only for the tables below.
MNIST digits are a far out-of-distribution set~\cite{lecun1998gradient}: same resolution, different pattern family.
Both datasets are normalized with the Fashion-MNIST training mean $0.286$ and standard deviation $0.353$.

\paragraph{Protocol.}
Both models share the CNN, the augmentation (horizontal flip and a random crop after 2-pixel padding), AdamW at learning rate $10^{-3}$, cosine decay, batch size 128, and 12 epochs.
Backbone weight decay is $10^{-4}$.
Runs use two seeds, $\{0,1\}$, on Apple MPS (PyTorch 2.14).
Validation probabilities for D-VBLL use 16 logit samples; the test tables use 32.
Reported means are over the two seeds.
Calibration is expected calibration error with 15 bins~\cite{guo2017calibration}.

\begin{table}[H]
\centering
\caption{Clean Fashion-MNIST test set. Mean over two seeds.}
\label{tab:clean}
\begin{tabular}{lcccc}
\toprule
Model & Accuracy & NLL & ECE & Entropy \\
\midrule
Softmax MAP & 0.930 & 0.200 & 0.0109 & 0.168 \\
D-VBLL ($1/T$) & 0.928 & 0.200 & 0.0075 & 0.183 \\
\bottomrule
\end{tabular}
\end{table}

On clean clothes the two models are tied on accuracy and likelihood.
D-VBLL's ECE is lower on both seeds (MAP $0.0101$ and $0.0116$; D-VBLL $0.0075$ and $0.0076$).
The absolute gap is small because the MAP model is already well calibrated.
Figure~\ref{fig:curves} shows the shared learning curve.
Figure~\ref{fig:rel} shows that D-VBLL's reliability curve sits closer to the diagonal through the middle bins.

\begin{figure}[H]
\centering
\safeimg[0.95\linewidth]{curves.png}
\caption{Validation accuracy and NLL. Bands are the two seeds. The training losses are not plotted together because cross-entropy and the ELBO are different objectives.}
\label{fig:curves}
\end{figure}

\begin{figure}[H]
\centering
\safeimg[0.62\linewidth]{reliability.png}
\caption{Reliability diagram on the clean test set, seed 0. Both models are close to the diagonal; D-VBLL is the smoother of the two.}
\label{fig:rel}
\end{figure}

\paragraph{Where the errors are.}
Seed-0 D-VBLL per-class accuracy is high on bags ($0.989$), trousers ($0.987$), and sandals ($0.978$), and lowest on shirts ($0.769$).
The largest off-diagonal counts are shirt$\to$T-shirt (96), T-shirt$\to$shirt (79), shirt$\to$coat (70), and coat$\to$shirt (41).
Those pairs share a torso silhouette.
Ankle boot$\to$sneaker (40) is the main footwear confusion.
Figure~\ref{fig:cm} is the row-normalized matrix.
Figure~\ref{fig:gallery} shows the failure mode the ECE number hides: a few errors still receive predictive probability near 1, including a boot read as a sneaker.

\begin{figure}[H]
\centering
\safeimg[0.72\linewidth]{confusion.png}
\caption{D-VBLL confusion on the clean test set, seed 0. Mass off the diagonal sits in the shirt / T-shirt / coat / pullover block.}
\label{fig:cm}
\end{figure}

\begin{figure}[H]
\centering
\safeimg[0.95\linewidth]{gallery.png}
\caption{D-VBLL seed 0. Titles are the predictive class, its probability, and the true class. Red titles are mistakes, including high-confidence ones.}
\label{fig:gallery}
\end{figure}

\paragraph{Corruptions.}
Corruptions are applied in pixel space at test time only.
The models were not trained on them.
Table~\ref{tab:corr} and Figure~\ref{fig:corr} show that geometric and photometric shift hurts both heads.
At noise $0.30$, rotation $60^{\circ}$, and low contrast, accuracy falls to roughly $0.12$--$0.23$.
Low contrast is the instructive case: accuracy is about $0.19$--$0.23$, but mean entropy stays near $0.33$, lower than under noise, and ECE exceeds $0.65$.
Both models are confidently wrong when the silhouette is preserved and the intensity is flattened.
Last-layer weight uncertainty is not a corruption detector.

\begin{table}[H]
\centering
\caption{Test accuracy under corruption. Mean over two seeds.}
\label{tab:corr}
\begin{tabular}{lcc}
\toprule
Condition & Softmax MAP & D-VBLL \\
\midrule
Clean & 0.930 & 0.928 \\
Gaussian noise $0.15$ & 0.422 & 0.373 \\
Gaussian noise $0.30$ & 0.145 & 0.180 \\
Rotate $30^{\circ}$ & 0.276 & 0.287 \\
Rotate $60^{\circ}$ & 0.121 & 0.117 \\
$5\times 5$ average blur & 0.409 & 0.430 \\
Low contrast & 0.193 & 0.235 \\
\bottomrule
\end{tabular}
\end{table}

\begin{figure}[H]
\centering
\safeimg[0.98\linewidth]{corruption_bars.png}
\caption{Accuracy falls under every corruption. Entropy rises under noise and rotation, but stays low under low contrast, where both models are confidently wrong.}
\label{fig:corr}
\end{figure}

\paragraph{Out-of-distribution digits.}
Using maximum predictive probability as the in-distribution score, AUROC against MNIST is $0.902$ and $0.808$ for the two MAP seeds (mean $0.855$), and $0.825$ and $0.855$ for D-VBLL (mean $0.840$).
The seed spread is larger than the gap between models, so I do not claim an OOD win for either head.
Figure~\ref{fig:ood} shows why the AUROC is only moderate: most clothes sit near probability 1, and a visible slice of digits does too.

\begin{figure}[H]
\centering
\safeimg[0.95\linewidth]{ood_hist.png}
\caption{Maximum predictive probability, seed 0. MNIST is shifted left of Fashion-MNIST, but both distributions have a spike at 1.}
\label{fig:ood}
\end{figure}

\paragraph{Why the Bayesian head stays close to MAP.}
After training at KL weight $1/T$, the mean diagonal variance of $W$ is about $4.9\times 10^{-4}$.
The prior variance is $2/D=0.0156$, roughly thirty times larger, so the posterior has concentrated.
The KL itself is about $2.3\times 10^{3}$; divided by $T=54000$ it contributes only a few hundredths to the loss, which is enough to regularize but not enough to keep the prior's uncertainty.
Figure~\ref{fig:norm} then has a simple shape for both models: predictive entropy falls as $\|\phi\|$ grows.
Large features produce large mean logits, and the quadratic variance term is too small to cancel that peak.
This is the opposite of a distance-aware uncertainty method, and it matches the modest ECE gap in Table~\ref{tab:clean}.

\begin{figure}[H]
\centering
\safeimg[0.72\linewidth]{norm_entropy.png}
\caption{Clean-test entropy against feature norm. Higher-norm features are more confident for both heads.}
\label{fig:norm}
\end{figure}

\paragraph{KL-weight ablation (seed 0).}
The $1/T$ factor is a hyperparameter of the ELBO estimator, not a constant of nature.
I retrained D-VBLL at $10/T$ and $100/T$ with the same seed and schedule.
Table~\ref{tab:kl} shows the trade.
Stronger KL pulls $q(W)$ toward the prior.
Mean weight variance goes from $4.9\times 10^{-4}$ at $1/T$ to $4.7\times 10^{-3}$ at $10/T$ and $8.1\times 10^{-3}$ at $100/T$, against a prior variance of $0.0156$.
The KL value falls from about $2.3\times 10^{3}$ to $229$, and predictive entropy rises from $0.187$ to $0.639$.
Clean accuracy stays at $0.927$--$0.929$ across all three settings.
What changes is the shape of the predictive distribution: clean NLL and ECE get worse, because these test images want sharp probabilities, while MNIST AUROC on this seed rises from $0.825$ to $0.916$.
$1/T$, the setting in the main table, is the best calibrated of the three.
The ablation shows that the KL term is alive: it trades clean calibration for broader uncertainty, without a collapse in accuracy at this model size.

\begin{table}[H]
\centering
\caption{D-VBLL KL multiplier, seed 0, clean test unless noted. Weight variance is the mean diagonal entry of $q(W)$.}
\label{tab:kl}
\begin{tabular}{lccccc}
\toprule
KL weight & Accuracy & NLL & ECE & Entropy & MNIST AUROC \\
\midrule
$1/T$ & 0.929 & 0.201 & 0.0075 & 0.187 & 0.825 \\
$10/T$ & 0.927 & 0.212 & 0.0154 & 0.280 & 0.869 \\
$100/T$ & 0.927 & 0.289 & 0.0937 & 0.639 & 0.916 \\
\bottomrule
\end{tabular}
\end{table}

\begin{figure}[H]
\centering
\safeimg[0.95\linewidth]{kl_ablation.png}
\caption{Seed-0 ablation of the KL multiplier. $1/T$ is the main setting.}
\label{fig:kl}
\end{figure}

\section{What I Have Learnt from This AI Assignment}
\begin{itemize}
  \item The useful constraint was the brief itself: probability or Bayes, recent, and not a repeat of FasterNet. Once that was stated, the agent could reject a fast CNN and look for a last-layer posterior with a checkable formula.
  \item A Bayesian classifier is defined by the predictive distribution, not by the word ``uncertainty'' in the README. Here that distribution is an average of softmaxes under $q(W)$. The self-check on Equation~\eqref{eq:bound} mattered more than the first epoch's accuracy.
  \item Matching accuracy is a real result. On this CNN, D-VBLL at $1/T$ ties MAP on accuracy and NLL and wins only on ECE, by a small margin that shows up on both seeds. I should not describe that as a large robustness gain.
  \item Posterior variance can collapse. The prior said $0.016$; training at $1/T$ produced about $5\times 10^{-4}$. After that, sampling the last layer cannot invent distance-aware uncertainty. Raising the KL weight reverses the collapse: entropy and MNIST AUROC go up, clean ECE gets worse, and accuracy on this CNN barely moves.
  \item Confidence is not a corruption alarm. Low contrast keeps the garment shape, both models stay confident, and ECE exceeds $0.65$. Shirt versus T-shirt is a label-overlap problem, not a Bayes problem.
  \item Out-of-distribution AUROC needs more than one seed. MAP's two seeds differed by about nine points. A single-run bar chart would have picked a winner that the second seed does not support.
\end{itemize}

\section{Source Code Webpage Link}
Source code:\\[0.3em]
\url{https://github.com/AngusJai/CISC3024-AI-Assignment2}\\[0.3em]
{\small The repository contains \texttt{models/}, \texttt{train.py}, \texttt{evaluate.py}, \texttt{run\_all.py}, \texttt{run\_ablation.py}, \texttt{process\_log.md}, \texttt{report.tex}, and the metrics and figures under \texttt{outputs/}. Dataset downloads (\texttt{data/}) and virtual environments are gitignored. Trained checkpoints for the main seeds are included so \texttt{demo.py} can be run without retraining.}

\section{Conclusion}
D-VBLL is a recent Bayesian classifier that an agent can implement from a variational bound and test on clothing images in one afternoon.
With the paper's $1/T$ KL weight it reproduces the softmax baseline's accuracy and slightly improves calibration.
It does not automatically detect digit images or low-contrast clothes, because the approximate posterior concentrates and the predictive distribution collapses toward the mean network.
The KL ablation makes that statement empirical: turn the KL up and entropy and digit-detection AUROC rise, while clean calibration falls. Accuracy on this network stays near $0.93$.

\begin{thebibliography}{9}
\bibitem{harrison2024vbll}
J.~Harrison, J.~Willes, and J.~Snoek.
Variational Bayesian Last Layers.
ICLR, 2024.
\bibitem{xiao2017fashion}
H.~Xiao, K.~Rasul, and R.~Vollgraf.
Fashion-MNIST: a Novel Image Dataset for Benchmarking Machine Learning Algorithms.
arXiv:1708.07747, 2017.
\bibitem{lecun1998gradient}
Y.~LeCun, L.~Bottou, Y.~Bengio, and P.~Haffner.
Gradient-based learning applied to document recognition.
Proceedings of the IEEE, 1998.
\bibitem{hendrycks2017baseline}
D.~Hendrycks and K.~Gimpel.
A Baseline for Detecting Misclassified and Out-of-Distribution Examples in Neural Networks.
ICLR, 2017.
\bibitem{guo2017calibration}
C.~Guo, G.~Pleiss, Y.~Sun, and K.~Q.~Weinberger.
On Calibration of Modern Neural Networks.
ICML, 2017.
\bibitem{blundell2015weight}
C.~Blundell, J.~Cornebise, K.~Kavukcuoglu, and D.~Wierstra.
Weight Uncertainty in Neural Networks.
ICML, 2015.
\bibitem{hoffman2013svi}
M.~D.~Hoffman, D.~M.~Blei, C.~Wang, and J.~Paisley.
Stochastic Variational Inference.
Journal of Machine Learning Research, 2013.
\end{thebibliography}

\end{document}
